% 证明需要两项前提；连接表示条件汇合，不表示一致收敛产生经验最优性。
\begin{tikzpicture}[figure picture,foundation picture,x=1mm,y=1mm]
 \path[use as bounding box] (0,-15) rectangle (169,28);
 \node[foundation heading,anchor=west] at (1,22) {空间层面的控制};
 \node[foundation input,fill=FigureBlueFill,minimum width=30mm] (u) at (16,8)
  {整个空间\\一致收敛};
 \node[foundation model,minimum width=30mm] (e) at (16,-8)
  {经验风险\\接近最优};
 \node[foundation operator] (j) at (42,0) {$+$};
 \node[foundation output,fill=FigureRedFill,minimum width=34mm] (r) at (63,0)
  {期望风险\\接近最优};
 \draw[foundation flow] (u.east)--(35,8)--(35,0)--(j);
 \draw[foundation flow] (e.east)--(35,-8)--(35,0);
 \draw[foundation flow] (j)--(r);

 \node[foundation heading,anchor=west] at (89,22) {算法层面的控制};
 \node[foundation input,fill=FigureBlueFill,minimum width=30mm] (s) at (104,8)
  {所选稳定性\\代价项$\beta_n$};
 \node[foundation model,minimum width=30mm] (f) at (104,-8)
  {经验比较误差\\$\rho_n$};
 \node[foundation operator] (k) at (131,0) {$+$};
 \node[foundation output,fill=FigureRedFill,minimum width=34mm] (o) at (151,0)
  {平均超额风险\\$\leq\beta_n+\rho_n$};
 \draw[foundation flow] (s.east)--(124,8)--(124,0)--(k);
 \draw[foundation flow] (f.east)--(124,-8)--(124,0);
 \draw[foundation flow] (k)--(o);
\end{tikzpicture}
